\section*{अध्याय \ref{ch:b3:pericyclic} --- परिचक्रीय अभिक्रियाएँ}

\begin{solution}{exo:b3:pericyclic:1}
डील्स--ऐल्डर: [4+2] \omterm{def:b3:pericyclic:families}{चक्रीयोजन}। साइक्लोब्यूटीन विवरण: विद्युतचक्रीय। कोप: [3,3]
सिग्माट्रॉपी। सल्फ़ोलीन से $\ce{SO2}$ का निकास: कीलेट्रॉपी। ओज़ोन और ऐल्कीन: 1,3-द्विध्रुवी
([3+2]) \omterm{def:b3:pericyclic:families}{चक्रीयोजन}। $[1,5]$-H विस्थापन: सिग्माट्रॉपी।
\end{solution}

\begin{solution}{exo:b3:pericyclic:2}
हेक्साट्राईन (6 इलेक्ट्रॉन): तापीय \omterm{def:b3:pericyclic:rotation-modes}{विघूर्णी}, प्रकाशरासायनिक \omterm{def:b3:pericyclic:rotation-modes}{सहघूर्णी}।
ब्यूटाडाईन (4): तापीय \omterm{def:b3:pericyclic:rotation-modes}{सहघूर्णी}, प्रकाशरासायनिक \omterm{def:b3:pericyclic:rotation-modes}{विघूर्णी}।
\end{solution}

\begin{solution}{exo:b3:pericyclic:3}
चार इलेक्ट्रॉन, तापीय: \omterm{def:b3:pericyclic:rotation-modes}{सहघूर्णी} विवरण। दोनों मेथिल समूह, जो वलय के विपरीत
फलकों पर हैं, दोनों बाहर की ओर घूमते हैं: (\textit{E},\textit{E})-हेक्सा-2,4-डाईन।
\end{solution}

\begin{solution}{exo:b3:pericyclic:4}
\omterm{def:b3:pericyclic:faciality}{समफलकीय} [2+2] अभिक्रिया में $4q$ इलेक्ट्रॉन हैं: मूल अवस्था में एक
एथीन का HOMO और दूसरे का LUMO एक आबंधी और एक प्रतिआबंधी सिरे के साथ अतिव्यापन करते हैं।
उत्तेजित अवस्था में एक अणु का एकल-अधिकृत $\pi^*$
HOMO की भूमिका निभाता है, और दोनों सिरे मेल खाते हैं।
\end{solution}

\begin{solution}{exo:b3:pericyclic:5}
\omterm{def:b3:point-groups:mirror}{दर्पण तल} बना रहता है। हेक्साट्राईन: $\psi_1$ S, $\psi_2$ A, $\psi_3$ S अधिकृत; $\psi_4$ A, $\psi_5$ S, $\psi_6$ A रिक्त।
साइक्लोहेक्साडाईन, क्रम में: $\sigma$ S, $\pi_1$ S, $\pi_2$ A अधिकृत; $\pi_3^*$ S, $\pi_4^*$ A, $\sigma^*$ A रिक्त। क्रम में
सममिति से जोड़ने पर, $\psi_1 \to \sigma$, $\psi_2 \to \pi_2$, $\psi_3 \to \pi_1$: अधिकृत से अधिकृत, \omterm{def:b3:pericyclic:rotation-modes}{विघूर्णी} समापन
तापीय रूप से अनुमत है।
\end{solution}

\begin{solution}{exo:b3:pericyclic:6}\emergencystretch=3em
प्रकाश में छह इलेक्ट्रॉन: \omterm{def:b3:pericyclic:rotation-modes}{सहघूर्णी}, अतः दोनों अंतस्थ मेथिल समूह
विपरीत फलकों पर समाप्त होते हैं: \textit{trans}-5,6-डाइमेथिलसाइक्लोहेक्सा-1,3-डाईन (तापीय अभिक्रिया
\textit{cis} समावयवी देती है)।
\end{solution}

\begin{solution}{exo:b3:pericyclic:7}
C5 पर स्थित हाइड्रोजन डाईन के C1 पर, जो वलय के चारों ओर उसका पड़ोसी है, खिसकता है,
छह-इलेक्ट्रॉन संक्रमण अवस्था ($\sigma2_s + \pi4_s$) से होकर, \omterm{def:b3:pericyclic:faciality}{समफलकीय} रूप से: एक प्रवास
1-मेथिल- देता है, दूसरा 2-मेथिलसाइक्लोपेन्टाडाईन। पाँच-कार्बन शृंखला
हाइड्रोजन को एक फलक पर रहने देती है, जिसके लिए छोटा वलय वैसे भी विवश करता है।
\end{solution}

\begin{solution}{exo:b3:pericyclic:8}
$\pi2_s$ (ऐल्कीन) $+ \pi2_a$ (कीटीन C=C, जिस पर लंबवत C=O $\pi^*$ के माध्यम से विपरीत फलकों पर
आक्रमण होता है): $(4q + 2)_s$ घटकों की गणना एक है (ऐल्कीन), विषम:
तापीय रूप से अनुमत।
\end{solution}

\begin{solution}{exo:b3:pericyclic:9}
कुर्सी में वाइनिल समूह का C=C (परमाणु 1--2), ऑक्सीजन (3), $\ce{OCH2}$
कार्बन (4) और ऐलिल C=C (5--6); O--C आबंध टूटता है और C1--C6 आबंध बनता है:
उत्पाद पेन्ट-4-ईनैल है,
$\ce{OHC-CH2-CH2-CH=CH2}$।
\end{solution}

\begin{solution}{exo:b3:pericyclic:10}
आठ इलेक्ट्रॉन $4q$ हैं; एक \omterm{def:b3:pericyclic:faciality}{विपरीतफलकीय} घटक के साथ विन्यास मोबियस
सांस्थिति का है, जो $4q$ इलेक्ट्रॉनों के साथ ऐरोमैटिक है: संक्रमण अवस्था ऐरोमैटिक है और
अभिक्रिया तापीय रूप से अनुमत (जैसे ऑक्टाटेट्राईन का \omterm{def:b3:pericyclic:rotation-modes}{सहघूर्णी} समापन)।
\end{solution}

\begin{solution}{exo:b3:pericyclic:11}
सबसे बड़े गुणांकों को युग्मित करने से N3 अंतस्थ कार्बन से जुड़ता है:
1,4-द्विप्रतिस्थापित ट्राइऐज़ोल। उत्प्रेरक के बिना गुणांकों में थोड़ा ही अंतर होता है और
1,4 तथा 1,5 दोनों समावयवी बनते हैं; कॉपर(I) अभिक्रिया को एक
कॉपर ऐसीटिलाइड के माध्यम से चरणबद्ध बना देता है और केवल 1,4 समावयवी देता है।
\end{solution}

\begin{solution}{exo:b3:pericyclic:12}
हाइड्रोजन $j$-परमाणु मूलक के SOMO, $k = (j + 1)/2$, के सिरों को जोड़ता है। इसके अंतस्थ
गुणांक $\sin(k\pi/(j + 1))$ और $(-1)^{k+1}\sin(k\pi/(j + 1))$ के समानुपाती हैं। $j = 7$, $k = 4$ के लिए: विपरीत चिह्न, अतः
समान चिह्न की पालियाँ दोनों सिरों पर विपरीत फलकों पर हैं: \omterm{def:b3:pericyclic:faciality}{विपरीतफलकीय}।
$j = 5$, $k = 3$ के लिए: एक ही फलक पर समान चिह्न: \omterm{def:b3:pericyclic:faciality}{समफलकीय}।
\end{solution}

\begin{solution}{pb:b3:pericyclic:1}
\textbf{1.} वलय B का C9--C10 $\sigma$ आबंध; 5,7-डाईन के साथ यह एक हेक्साट्राईन देता है,
C10=C5--C6=C7--C8=C9।

\textbf{2.} छह इलेक्ट्रॉन (दोनों $\pi$ आबंध और $\sigma$ आबंध): एक विद्युतचक्रीय वलय-विवरण।

\textbf{3.} तापीय रूप से कोई साइक्लोहेक्साडाईन और उसका खुला हेक्साट्राईन केवल
धीरे, एक ऊँचे अवरोध से होकर, साम्यित होते हैं, और वलय अधिक स्थायी पक्ष है;
डाईन पराबैंगनी प्रकाश अवशोषित करता है, और उसकी उत्तेजित अवस्था
पिकोसेकंडों में खुल जाती है।

\textbf{4.} \omterm{def:b3:pericyclic:rotation-modes}{सहघूर्णी}: उत्तेजित अवस्था में ट्राईन
निकाय के एकल-अधिकृत LUMO ($k = 4$) के सिरों के चिह्न विपरीत हैं।

\textbf{5.} C6=C7 आबंध वलय से आता है: इसकी दोनों शृंखला-निरंतरताएँ, C5 और C8,
इसके एक ही ओर थीं, और वैसी ही रहती हैं: \textit{Z}।

\textbf{6.} तापीय विवरण, जो केवल \omterm{def:b3:pericyclic:rotation-modes}{विघूर्णी} रूप से अनुमत है, का अवरोध शरीर की ऊष्मा की क्षमता से
कहीं अधिक है, और यह चढ़ाई की ओर, कम स्थायी ट्राईन की ओर, ले जाता।

\textbf{7.} C19 (C10 पर स्थित मेथिल) से C9 तक: एक $[1,7]$-H सिग्माट्रॉपी विस्थापन।

\textbf{8.} आठ: ट्राईन के छह $\pi$ इलेक्ट्रॉन और C--H आबंध के दो।

\textbf{9.} $\sigma2 + \pi6$। \omterm{def:b3:pericyclic:faciality}{समफलकीय} रूप से दोनों s होते, प्रत्येक से एक $(4q + 2)_s$ घटक:
सम, निषिद्ध। ट्राईन के \omterm{def:b3:pericyclic:faciality}{विपरीतफलकीय} होने पर, $\sigma2_s + \pi6_a$: एक गिना गया घटक, विषम,
अनुमत।

\textbf{10.} कुंडलित, \textit{Z}-विन्यासी शृंखला में सात परमाणु हाइड्रोजन को एक फलक से
दूसरे पर जाने देते हैं; तीन-परमाणु शृंखला इतनी ऐंठ नहीं सकती।

\textbf{11.} मोबियस (एक \omterm{def:b3:pericyclic:faciality}{विपरीतफलकीय} घटक) आठ इलेक्ट्रॉनों, $4q$, के साथ: ऐरोमैटिक।

\textbf{12.} यह तापीय रूप से अनुमत है; शरीर के ताप पर इसका अवरोध धीरे पार होता है।

\textbf{13.} प्रकाश में छह इलेक्ट्रॉन: फिर \omterm{def:b3:pericyclic:rotation-modes}{सहघूर्णी}, परंतु यह सिरों को
दूसरी ओर घुमाकर H9 और C19 मेथिल को अन्य फलकों पर रख सकता है: एक वलय-बंद
त्रिविम समावयवी, ल्यूमिस्टेरॉल।

\textbf{14.} केंद्रीय आबंध \textit{E} होने पर C19 और C9 शृंखला के विपरीत ओर और बहुत
दूर होते हैं: हाइड्रोजन नहीं पहुँच सकता।

\textbf{15.} प्रीविटामिन D$_3$ स्वयं अवशोषित करता है और ल्यूमिस्टेरॉल और टैकिस्टेरॉल में बदल जाता है,
जो विटामिन D नहीं देते; एक प्रकाशस्थायी मिश्रण प्राप्त होता है, जो
उत्पादन को सीमित कर देता है।

\textbf{16.} विवरण उत्तेजित अवस्था में, फ़ेम्टोसेकंडों से पिकोसेकंडों में, होता है;
विस्थापन \qty{85}{kJ/mol} अवरोध वाली एक तापीय अभिक्रिया है।

\textbf{17.} $t_{1/2} = \ln 2/k = 0.693/\qty{0.023}{h^{-1}} = \qty{30}{h}$।

\textbf{18.} $1 - \eu^{-0.023 \times 24} = 0.42$।

\textbf{19.} $\ln 10/k = 2.303/\qty{0.023}{h^{-1}} = \qty{100}{h}$।

\textbf{20.} $k(\qty{293.15}{K}) = k(\qty{310.15}{K})\,\eu^{-(E_a/R)(1/293.15 - 1/310.15)} = 0.023 \times \eu^{-1.91} = \qty{3.4e-3}{h^{-1}}$।

\textbf{21.} $2.303/\qty{3.4e-3}{h^{-1}} = \qty{680}{h}$, लगभग 28 दिन।

\textbf{22.} तापीय चरण \qty{37}{\celsius} पर \qty{20}{\celsius} की तुलना में सात गुना तेज़ है, अतः
धूप में एक सुबह बना प्रीविटामिन कुछ ही दिनों में रूपांतरित हो जाता है।

\textbf{23.} अपने अनेक त्रिविमजनक केंद्रों सहित स्टेरॉयड ढाँचा स्टेरॉल से
बना-बनाया आता है; फिर प्रकाश और ऊष्मा दोनों परिचक्रीय चरण ठीक उसी तरह पूरे करते हैं जैसे
त्वचा में।

\textbf{24.} शरीर के ताप पर 90\,\% प्रीविटामिन D$_3$ लगभग
\qty{100}{h}, अर्थात् चार दिन, में विटामिन D$_3$ बन जाता है।
\end{solution}
